\section{Results}
\label{sec:results}

Every holdout measurement in this section is produced from the artifacts by the generator
described in Appendix~\ref{app:repro} and emitted into this document mechanically. The figures those
measurements are read against are not measurements of this holdout: registered expectations,
earlier cycles' results, and structural counts. Appendix~\ref{app:repro} records where each comes
from. The section answers RQ1 in
\S\ref{sec:results-primary}, RQ3 in \S\ref{sec:results-attribution}, and RQ2 in
\S\ref{sec:results-partners} through \S\ref{sec:results-adversarial}, before reporting the
evaluation controls and the computational cost.

\subsection{Baseline comparisons}
\label{sec:results-primary}

Against \fcheap{}, the registered primary baseline, the recorded implementation of \sestina{}
produced an observed win-rate edge of \vsevenBtwoFcheap{}~pp relative to parity over
\vsevenBtwoFcheapN{} games, with a nominal pooled 95\% interval of
$[\vsevenBtwoFcheapLo, \vsevenBtwoFcheapHi]$; this edge corresponds to a \vsevenBtwoFcheapWinRate{}
head-to-head win rate against \fcheap{}, not to a change from a separately measured baseline win
rate. The interval is formed by normal approximation from the per-bank deal-block bootstrap
intervals (\S\ref{sec:evaluation-intervals}), constructions that treat deals as independent
resampling units; the scored implementation retains cross-deal agent state, so independence and
coverage have not been established for this comparison (\S\ref{sec:discussion-engineering}). The
two bank estimates were \vsevenBtwoFcheapBankOne{} and
\vsevenBtwoFcheapBankTwo{}~pp. The pooled nominal lower bound exceeded the preregistered
\vsevenPrimaryFloor{}~pp threshold, both bank estimates were positive, and the configuration passed
the specified behavioural gate in its designated test mode (\S\ref{sec:design-gate}). The
registered primary decision rule was therefore satisfied on the nominal interval, and the
comparison \vsevenPrimaryVerdict{}; that verdict is a recorded protocol outcome rather than a
coverage-validated inference. The scored-mode isolation limitation is reported in
\S\ref{sec:results-controls}. Figure~\ref{fig:forest} places this comparison beside the others, and
Table~\ref{tab:frontier} gives the cells behind it.

\begin{figure}[!ht]\centering
{\tabfont\begin{tikzpicture}[x=1cm,y=1cm]
  \node[anchor=south west,font=\scriptsize\itshape,text=fishgray] at (0,0.300) {A. direct comparisons: SESTINA's edge against the named configuration};
  \node[anchor=south west,font=\scriptsize\itshape,text=fishgray] at (0,-3.930) {B. other estimands, oriented so that positive favours SESTINA};
  \node[anchor=east,font=\small] at (-0.25,0.000) {vs.\ v0.6 (deployed)};
  \draw[fishgray,thick] (5.789,0.000) -- (6.663,0.000);
  \draw[fishgray] (5.789,-0.110) -- (5.789,0.110);
  \draw[fishgray] (6.663,-0.110) -- (6.663,0.110);
  \fill[fishgray] (6.226,0.000) circle (0.075);
  \node[anchor=west,font=\scriptsize\ttfamily] at (8.250,0.000) {+4.63 [+4.19, +5.06]};
  \node[anchor=east,font=\small] at (-0.25,-0.720) {vs.\ F-cheap \textit{(primary)}};
  \draw[fishgreen,ultra thick] (4.491,-0.720) -- (5.383,-0.720);
  \draw[fishgreen] (4.491,-0.830) -- (4.491,-0.610);
  \draw[fishgreen] (5.383,-0.830) -- (5.383,-0.610);
  \fill[fishgreen] (4.937,-0.720) circle (0.075);
  \node[anchor=west,font=\scriptsize\ttfamily] at (8.250,-0.720) {+3.33 [+2.88, +3.78]};
  \node[anchor=east,font=\small] at (-0.25,-1.440) {vs.\ F-mid};
  \draw[fishgray,thick] (3.621,-1.440) -- (5.388,-1.440);
  \draw[fishgray] (3.621,-1.550) -- (3.621,-1.330);
  \draw[fishgray] (5.388,-1.550) -- (5.388,-1.330);
  \fill[fishgray] (4.505,-1.440) circle (0.075);
  \node[anchor=west,font=\scriptsize\ttfamily] at (8.250,-1.440) {+2.89 [+2.00, +3.78]};
  \node[anchor=east,font=\small] at (-0.25,-2.160) {vs.\ phase-2 composite};
  \draw[fishgray,thick] (1.340,-2.160) -- (2.215,-2.160);
  \draw[fishgray] (1.340,-2.270) -- (1.340,-2.050);
  \draw[fishgray] (2.215,-2.270) -- (2.215,-2.050);
  \fill[fishgray] (1.777,-2.160) circle (0.075);
  \node[anchor=west,font=\scriptsize\ttfamily] at (8.250,-2.160) {+0.15 [-0.29, +0.59]};
  \node[anchor=east,font=\small] at (-0.25,-2.880) {vs.\ v0.5};
  \draw[fishgray,thick] (6.157,-2.880) -- (7.404,-2.880);
  \draw[fishgray] (6.157,-2.990) -- (6.157,-2.770);
  \draw[fishgray] (7.404,-2.990) -- (7.404,-2.770);
  \fill[fishgray] (6.781,-2.880) circle (0.075);
  \node[anchor=west,font=\scriptsize\ttfamily] at (8.250,-2.880) {+5.18 [+4.56, +5.81]};
  \node[anchor=east,font=\small] at (-0.25,-4.350) {worst partner substitution};
  \draw[fishgray,thick] (0.596,-4.350) -- (2.297,-4.350);
  \draw[fishgray] (0.596,-4.460) -- (0.596,-4.240);
  \draw[fishgray] (2.297,-4.460) -- (2.297,-4.240);
  \fill[fishgray] (1.446,-4.350) circle (0.075);
  \node[anchor=west,font=\scriptsize\ttfamily] at (8.250,-4.350) {-0.19 [-1.04, +0.67]};
  \node[anchor=east,font=\small] at (-0.25,-5.070) {vs.\ strongest adversary};
  \draw[fishgray,thick] (3.883,-5.070) -- (5.159,-5.070);
  \draw[fishgray] (3.883,-5.180) -- (3.883,-4.960);
  \draw[fishgray] (5.159,-5.180) -- (5.159,-4.960);
  \fill[fishgray] (4.521,-5.070) circle (0.075);
  \node[anchor=west,font=\scriptsize\ttfamily] at (8.250,-5.070) {+2.91 [+2.27, +3.55]};
  \draw[fishgray,thin] (-4.2,-3.430) -- (11.400,-3.430);
  \draw[densely dashed,fishgray] (1.633,-5.620) -- (1.633,0.620);
  \draw[->,fishgray] (0,-5.620) -- (8.100,-5.620);
  \draw[fishgray] (0.639,-5.620) -- (0.639,-5.720) node[below,font=\scriptsize] {-1};
  \draw[fishgray] (1.633,-5.620) -- (1.633,-5.720) node[below,font=\scriptsize] {0};
  \draw[fishgray] (2.626,-5.620) -- (2.626,-5.720) node[below,font=\scriptsize] {1};
  \draw[fishgray] (3.619,-5.620) -- (3.619,-5.720) node[below,font=\scriptsize] {2};
  \draw[fishgray] (4.612,-5.620) -- (4.612,-5.720) node[below,font=\scriptsize] {3};
  \draw[fishgray] (5.605,-5.620) -- (5.605,-5.720) node[below,font=\scriptsize] {4};
  \draw[fishgray] (6.599,-5.620) -- (6.599,-5.720) node[below,font=\scriptsize] {5};
  \draw[fishgray] (7.592,-5.620) -- (7.592,-5.720) node[below,font=\scriptsize] {6};
  \node[font=\scriptsize,anchor=north] at (4.000,-6.070) {win-rate edge (percentage points); positive favours SESTINA};
\end{tikzpicture}
}
\caption{Principal comparisons for \sestina{}. Each point is a pooled win-rate edge in percentage
points over both holdout banks (\S\ref{sec:evaluation-intervals}: cross-bank pooling of per-bank
deal-clustered bootstrap intervals; the bars are nominal 95\% intervals). Both panels are oriented so that
positive values favour \sestina{}. The upper panel shows the five direct head-to-head comparisons,
with the registered primary comparison highlighted. The lower panel shows two different estimands:
the worst partner-substitution row of \S\ref{sec:results-partners}, which is a difference between
two upgrade effects (\sestina{}'s upgrade over \vsix{} minus \vsix{}'s upgrade over \vfive{}, with
the named partner) rather than a head-to-head edge, and \sestina{}'s edge against the strongest of
the \vsevenBfourArms{} adversarial responses of \S\ref{sec:results-adversarial}, shown here as the
negation of the adversary's edge printed in Table~\ref{tab:adversaries}, whose convention is the
adversary's edge against \sestina{}.}
\label{fig:forest}
\end{figure}

\begin{table}[!ht]\centering
\caption{Strength against the \vsix{} frontier and lineage baselines. Each row is \sestina{}'s
win-rate edge in percentage points against the named configuration, paired duplicate blocks, pooled
over both holdout banks with the per-bank estimates beside it; intervals are nominal 95\%
intervals as constructed in \S\ref{sec:evaluation-intervals}.}
\label{tab:frontier}
{\tabfont\input{tables_v07/frontier}}
\end{table}

Against the lineage baselines, \sestina{} gains \vsevenBtwoIncumbent{}~pp over the deployed
\vsix{} policy $[\vsevenBtwoIncumbentLo, \vsevenBtwoIncumbentHi]$, \vsevenBtwoFmid{}~pp over \fmid{}
$[\vsevenBtwoFmidLo, \vsevenBtwoFmidHi]$, and \vsevenBtwoVfive{}~pp over v0.5.

Against \composite{} it gains \vsevenBtwoComposite{}~pp
$[\vsevenBtwoCompositeLo, \vsevenBtwoCompositeHi]$, an interval containing zero.
\S\ref{sec:results-composite} treats that result, which the protocol had named in advance as a
condition of non-confirmation.

\paragraph{Descriptive declaration accuracy.} The declaration-family floor applies to
the declaration-accuracy column. The harness emits declaration accuracy as a bare rate with no
interval of any kind, so no interval can be quoted and no claim about the declaration channel is
made. \sestina{} runs \vsevenDeclFrozenLo{} to \vsevenDeclFrozenHi{} across the five cells against
opponents' \vsevenDeclOppLo{} to \vsevenDeclOppHi{}. That is a diagnostic.

\subsection{Secondary criteria and diagnostic comparisons}
\label{sec:results-nonconfirm}

The protocol prespecified seven non-confirmation conditions. Two arose, and they are reported here
rather than in a discussion of caveats; the remaining diagnostic comparisons follow.

\subsubsection{Phase-2 composite comparison}
\label{sec:results-composite}

\sestina{}'s comparison with \composite{} was \vsevenBtwoComposite{}~pp (nominal 95\% CI
$[\vsevenBtwoCompositeLo, \vsevenBtwoCompositeHi]$) over \vsevenBtwoCompositeN{} games. Both banks
are positive, so the sign replicates, but the interval includes zero, so the evaluation does not
resolve an additional benefit over that composite.

The registered interpretation is explicit, and this paper adopts it. The increment that an earlier
phase had reported over the composite, \vsevenKthreeOverComposite{}
$[\vsevenKthreeOverCompositeLo, \vsevenKthreeOverCompositeHi]$ on a single bank, did not replicate
on holdout.

The arithmetic makes the cell interpretable. \sestina{} minus \composite{} is exactly the tie-break
repair, the urgency removal and the stall detector added, with one further switch removed.
\S\ref{sec:results-attribution} measures the stall detector's contribution as identically zero, and
the removed switch's leave-one-out drop was zero on training material. The protocol states that
group's training value in advance as \vsevenTrainGroup{} $[\vsevenTrainGroupLo, \vsevenTrainGroupHi]$;
on holdout it is \vsevenBtwoComposite{} $[\vsevenBtwoCompositeLo, \vsevenBtwoCompositeHi]$.

This result does not alter the positive comparison with \fcheap{}, which is a different baseline
(Table~\ref{tab:configs}): \fcheap{} is a configuration of the \vsix{} frontier that existed before
this cycle, and \sestina{} improves on it by a replicated margin. What the result does is limit the
attribution of further improvement to the modifications introduced after the composite was
assembled; it does not establish that those modifications contribute nothing.

\subsubsection{Component-attribution criterion}
\label{sec:results-location}

The protocol's third non-confirmation condition is that the lattice cannot locate the gain: if
removing every single component in turn costs less than a third of the whole, the configuration has a
gain it cannot explain. Read per bank, as \S\ref{sec:evaluation-replication} requires, the two banks
disagree (Table~\ref{tab:location}).

\begin{table}[!ht]\centering
\caption{The attribution location test, read per bank. \emph{Whole} is the measured whole,
\emph{one third} the protocol's threshold, and \emph{largest single drop} the biggest
leave-one-out drop the lattice contains.}
\label{tab:location}
{\tabfont
\begin{tabular}{lrrlc}
\toprule
Bank & Whole & One third & Largest single drop & Verdict \\
\midrule
7090001 & \vsevenLocWholeOne{} & \vsevenLocThirdOne{} & \texttt{\vsevenLocAtOne{}} \vsevenLocBiggestOne{} & \vsevenLocVerdictOne{} \\
7090003 & \vsevenLocWholeThree{} & \vsevenLocThirdThree{} & \texttt{\vsevenLocAtThree{}} \vsevenLocBiggestThree{} & \vsevenLocVerdictThree{} \\
\bottomrule
\end{tabular}}
\end{table}

The two banks' largest drops differ by \vsevenLocSplit{}~pp against a half-width on that difference of
comparable size, so the split is real on the point estimates and is not itself statistically resolved.
Under the replication criterion the pooled verdict is not replicated
(\vsevenLocReplicates{}), and the supported interpretation is that this battery does not settle
whether the gain is located by the protocol's one-third rule. Note that the disagreement is between verdict labels
rather than between signs, which is itself a recorded extension of the criterion. The sub-additivity
claim of \S\ref{sec:results-attribution} passes on both banks; it is the location test specifically
that splits.

\subsubsection{Selection-bias diagnostics}
\label{sec:results-selection}

\vsevenSelK{} configurations were scored against the deployed policy during development. At the
lattice cell size the per-cell standard deviation is \vsevenSelSigma{}~pp, so the quantity the protocol registers as the expected maximum of
\vsevenSelK{} draws under the null that none differs from the reference is
$\sigma\sqrt{2\ln K} = \vsevenSelExpectedMax{}$~pp, which is an upper bound on that
expectation rather than the expectation itself.

The held-out deals were not used for configuration selection, so the selection term on holdout
material is \vsevenSelOnHoldout{}~pp: this removes direct selection on those outcomes, subject to
the documented separation procedure of \S\ref{sec:evaluation-seal}, and it is a property of the
design rather than a measurement.
Table~\ref{tab:selection} sets the four registered quantities against their holdout measurements.

\begin{table}[!ht]\centering
\caption{The four quantities the registered selection check tabulates, against their holdout
measurements. Three reproduce or exceed; one falls short. These are not every figure the protocol
states in advance, and two of the partner-regime predictions in \S\ref{sec:results-partners} also
come in low.}
\label{tab:selection}
{\tabfont\input{tables_v07/selection}}
\end{table}

\vsevenSelReproduce{} of the four reproduce or exceed their training values. The largest shortfall,
\vsevenSelWorstShortfall{}~pp against the registered \vsevenSelExpectedMax{},
falls on the composite cell, which is precisely the quantity measuring the group of components the
configuration was selected over. That is the pattern the check was written to look for and it appears
where selection occurred. One quantity cannot establish that the group's training value was entirely
selection effect, and the check is not a formal correction
(\S\ref{sec:evaluation-multiplicity}); what it provides is a consistency test that the data pass in
the direction \S\ref{sec:results-composite} indicates.

\subsubsection{Preregistered non-confirmation criteria}
\label{sec:results-seven}

Table~\ref{tab:seven} answers each of the seven in the protocol's own order. The recorded verdicts
apply the protocol's rules as written; the third column states the reading each rule uses, because
the rules are not uniform: conditions 1 and 5 are decided on interval bounds, condition 2 on whether
the interval contains zero, condition 3 per bank on point estimates, and conditions 4 and 7 on point
estimates against thresholds that the deciding intervals cross.

\begin{table}[!ht]\centering
\caption{The seven registered non-confirmation conditions against what the holdout returned.
\emph{Recorded verdict} is the verdict under the protocol's rule as written; \emph{reading with
uncertainty} states the statistic that rule uses and, where the deciding interval crosses the
threshold, says so. Every interval is nominal (\S\ref{sec:evaluation-intervals}).
Appendix~\ref{app:thresholds} reproduces the criteria with the same distinction.}
\label{tab:seven}
{\tabfont
\begin{tabularx}{\linewidth}{@{}c Q{0.8} Q{1.1} Q{1.1}@{}}
\toprule
\# & Condition & Recorded verdict & Reading with uncertainty \\
\midrule
1 & the frontier comparison & not met. \vsevenBtwoFcheap{} $[\vsevenBtwoFcheapLo, \vsevenBtwoFcheapHi]$, sign replicates & passes the registered rule using nominal intervals (lower bound above \vsevenPrimaryFloor{}); coverage unestablished \\
2 & the composite comparison & met. \vsevenBtwoComposite{} $[\vsevenBtwoCompositeLo, \vsevenBtwoCompositeHi]$ & interval contains zero; the comparison is unresolved rather than negative \\
3 & attribution: the location test & does not replicate. \vsevenLocVerdictOne{} on one bank, \vsevenLocVerdictThree{} on the other & per-bank point estimates against one third of the whole; not resolved \\
4 & brittleness under partner change or cross-play & not met. Minimum \vsevenSoneMin{}; cross-play off-diagonal minus diagonal \vsevenXpGap{} & point estimates; the deciding partner interval $[\vsevenOppFiveWithholderLo, \vsevenOppFiveWithholderHi]$ crosses the $\vsevenSoneThreshold{}$ threshold, and the cross-play contrast carries no interval \\
5 & a successful adversary & not met. No search found a positive edge; largest upper bound \vsevenBfourWorstUpper{} & passes the registered rule using nominal intervals for the tested searches (upper bounds below the class floor); coverage unestablished \\
6 & the behavioural validity gate & not met. All \vsevenGateRules{} rules pass & deterministic checks on training material \\
7 & the validity controls & not met. Every control behaves as specified on the point-estimate reading & point estimates; on the lower-bound reading the \texttt{0.08} rung is not recovered and the sub-floor rung's upper bound exceeds the floor (\S\ref{sec:results-controls}) \\
\bottomrule
\end{tabularx}}
\end{table}

Two of the seven are met or unresolved, being condition 2 outright and condition 3 by
non-replication. The other five are not met on the recorded reading; conditions 4 and 7 are not
resolved against their thresholds once uncertainty is taken into account.

\subsection{Opponent-panel performance}
\label{sec:results-panel}

\vsevenPanelCells{} cells. Four configurations were scored against the same
\vsevenPanelMembers{}-member panel (\vsevenPanelDistinct{} distinct policies, \S\ref{sec:evaluation-panel})
on the same two banks: \vsevenPanelNear{} near-class members at \vsevenDealsNear{} deals per bank,
\vsevenPanelFar{} far-class archetypes at \vsevenDealsFar{}, and \vsevenPanelExpensive{} expensive
members at \vsevenDealsExpensive{}. All four carry a cell for every member, which is what makes the
empirical regret of \S\ref{sec:evaluation-panel} defined. Table~\ref{tab:worstcase} gives each
configuration's worst cell. The extrema and rankings in this subsection are computed from pooled
point estimates and are descriptive; the interval beside a worst cell is that cell's own interval.

\begin{table}[!ht]\centering
\caption{Worst panel cell per configuration: the panel member against which the configuration's
pooled win-rate edge (percentage points, both banks) is lowest, with that cell's nominal 95\%
interval.
Per-bank values are given in bank order, 7090001 then 7090002.}
\label{tab:worstcase}
{\tabfont\input{tables_v07/worstcase}}
\end{table}

\sestina{}'s worst cell is \vsevenWorstFrozen{}~pp
$[\vsevenWorstFrozenLo, \vsevenWorstFrozenHi]$ against the deployed policy's \vsevenWorstIncumbent{}
$[\vsevenWorstIncumbentLo, \vsevenWorstIncumbentHi]$, satisfying the registered secondary claim that
the worst case not be catastrophic.

Two qualifications belong with that number.

First, it does not replicate in sign. \sestina{}'s worst cell is
\vsevenWorstFrozenBankOne{} on the first bank and \vsevenWorstFrozenBankTwo{} on the second, attained
at a different panel member on each (\texttt{\vsevenWorstFrozenBankOneOpp{}} and
\texttt{\vsevenWorstFrozenBankTwoOpp{}}); on one bank \sestina{} has no negative cell anywhere on
the panel. Under the criterion of \S\ref{sec:evaluation-replication} the pooled value is not
replicated. What survives replication is the comparison the protocol actually gated on, since
\sestina{}'s worst cell is far above the deployed policy's on both banks, rather than the pooled
value itself.

Second, it is not the best worst case of the four. Configurations that are themselves panel
members meet themselves in a mirror cell that is exactly zero, which can only lower their worst
case (\S\ref{sec:evaluation-panel}), so the worst case is printed both ways; \composite{} is ahead
of \sestina{} on both readings, at \vsevenWorstComposite{} with its self-cell and
\vsevenWorstCompositeExSelf{} without.

\begin{table}[!ht]\centering
\caption{Empirical minimax regret over the shared panel (\S\ref{sec:evaluation-panel}), in
percentage points of win-rate edge on pooled point estimates, best first, with the panel member at
which each configuration's largest shortfall from the best of the four is attained. No interval is
attached to this statistic. The protocol sets no threshold on it and requires only that it be
reported.}
\label{tab:regret}
{\tabfont\input{tables_v07/regret}}
\end{table}

\sestina{} is \vsevenRegretFrozenRank{}rd of four on empirical minimax regret
(Table~\ref{tab:regret}), at \vsevenRegretFrozen{} against \vsevenRegretBest{}'s
\vsevenRegretFcheap{} and \composite{}'s \vsevenRegretComposite{}. The two statistics rank the
configurations differently: the configuration with the best worst case is not the configuration
with the best regret. The regret difference is concentrated in the far archetypes
(Table~\ref{tab:farpanel}), opponents every configuration already beats by between
\vsevenFarSpanLo{} and \vsevenFarSpanHi{}~pp, which is a reason to weight it less but not a reason
to discount it.

\begin{table}[!ht]\centering
\caption{The far archetypes, where the minimax-regret differences are concentrated.}
\label{tab:farpanel}
{\tabfont\input{tables_v07/farpanel}}
\end{table}

Against the hardest sealed adversary the ordering reverses: on
\texttt{\vsevenHardSealed{}}, \sestina{} scores \vsevenHardFrozen{}
$[\vsevenHardFrozenLo, \vsevenHardFrozenHi]$ where the deployed policy scores
\vsevenHardIncumbent{} and \fcheap{} scores \vsevenHardFcheap{}. \vsevenHardBeaters{} of the four
configurations beat that member with a lower bound above zero, namely \vsevenHardBeaterNames{}.

\paragraph{Descriptive panel means.} Over the \vsevenPanelMembers{} shared cells, which comprise
\vsevenPanelDistinct{} distinct policies since three are named twice and are therefore
double-counted by any mean, the figures are \sestina{} \vsevenPanelMeanFrozen{},
\composite{} \vsevenPanelMeanComposite{}, \fcheap{} \vsevenPanelMeanFcheap{} and the deployed policy
\vsevenPanelMeanIncumbent{}. The reporting rule of \S\ref{sec:evaluation-panel} forbids leading with
this, and it points the opposite way from minimax regret. Appendix~\ref{app:panel} gives every cell.

\subsection{Component ablations}
\label{sec:results-attribution}

The registered prediction was that the contestation weight of \S\ref{sec:agent-contest} carries the
gain, checkable as a leave-one-out drop of roughly two percentage points. On that check it holds:
the holdout drop is \vsevenLooRtwelve{}~pp $[\vsevenLooRtwelveLo, \vsevenLooRtwelveHi]$, with the
training figure \vsevenTrainRtwelveLoo{} inside the interval, and an add-one-in of
\vsevenAddInRtwelve{}. It is the largest observed leave-one-out drop of the five, which is a
descriptive ordering rather than a statistically established largest contribution, since the
drops' intervals overlap. The stronger statement, that this coordinate is where the gain lives, is
not supported at this precision: the drop's lower bound is \vsevenLooRtwelveLo{}, and
\S\ref{sec:results-location} reports the protocol's own location test failing to replicate.
Table~\ref{tab:attribution} is the lattice.

\begin{table}[!ht]\centering
\caption{The attribution lattice, in percentage points of win-rate edge against the deployed
\vsix{} policy on the two lattice banks. Add-one-in is the edge of the deployed policy plus the
named component; leave-one-out is the drop from \sestina{}'s edge when the component is removed, a
between-cell contrast whose interval is propagated as in \S\ref{sec:evaluation-intervals}.
Reference cell \sestina{} against the deployed policy on the lattice banks: \vsevenBfiveWhole{}
$[\vsevenBfiveWholeLo, \vsevenBfiveWholeHi]$. Harness identifiers are given in parentheses.}
\label{tab:attribution}
{\tabfont\begin{tabular}{lrcrc}
\toprule
Component & Add-one-in & 95\% CI & Leave-one-out & 95\% CI \\
\midrule
half-suit contestation (\texttt{r12=25}) & $+2.18$ & $[+1.55, +2.81]$ & $+1.69$ & $[+0.80, +2.57]$ \\
guarded search (\texttt{s1=1}) & $+1.91$ & $[+1.38, +2.44]$ & $+0.57$ & $[-0.33, +1.46]$ \\
public-history tie-break (\texttt{rtie=1}) & $+0.71$ & $[+0.08, +1.34]$ & $+1.08$ & $[+0.20, +1.97]$ \\
urgency escalation off (\texttt{urgency-off}) & $+1.37$ & $[+1.12, +1.63]$ & $+0.64$ & $[-0.25, +1.53]$ \\
stall detector (\texttt{stall=12}) & $+0.00$ & $[+0.00, +0.00]$ & $+0.00$ & $[-0.89, +0.89]$ \\
rescoring switch (\texttt{m2=0}; not in the freeze) & $+0.77$ & $[+0.66, +0.88]$ & --- & --- \\
\bottomrule
\end{tabular}
}
\end{table}

The stall detector behaves exactly as predicted. It is worth \vsevenAddInStall{}
$[\vsevenAddInStallLo, \vsevenAddInStallHi]$ added and \vsevenLooStall{} removed, bit-identical on
holdout in win rate and in events per game alike: the rule does not fire.

The naive sum of the five components \sestina{} carries is \vsevenBfiveSum{}~pp against a measured
whole of \vsevenBfiveWhole{}~pp, a sub-additivity ratio of \vsevenSubAdditivity{}
$[\vsevenSubAdditivityLo, \vsevenSubAdditivityHi]$, which contains the phase-2 composition figure of
\vsevenSubAddPhaseTwo{}. Quoting the sum would overstate the configuration by
\vsevenSubAdditivityOverstate{}~pp. The registered attribution claim, that components are
sub-additive and no single leave-one-out drop exceeds the whole, is satisfied.

One component was measured but is not carried by \sestina{}. The \texttt{m2} switch is worth
\vsevenAddInMtwo{} $[\vsevenAddInMtwoLo, \vsevenAddInMtwoHi]$ added to the blueprint, an interval
excluding zero and larger than the tie-break's add-one-in. Its zero leave-one-out drop was measured
on training material and is not re-measured here, because \sestina{} does not carry the switch and
the protocol directs that it not be added back; there is no leave-one-out cell for it and there could
not be.

\subsection{Partner substitution and single-seat upgrades}
\label{sec:results-partners}

\vsevenSonePositive{} of \vsevenSoneRows{} changed-partner rows against a v0.5 opponent are positive
(Table~\ref{tab:partners}). The minimum is
\vsevenSoneMin{} $[\vsevenOppFiveWithholderLo, \vsevenOppFiveWithholderHi]$, which is above the
registered $\vsevenSoneThreshold{}$ collapse threshold on the point estimate but \emph{not resolved
against it}, since the interval's lower bound falls below the threshold. This battery does not
separate a pass from a failure on the row that decides the claim, and reading a threshold from a
point estimate is a move this paper declines elsewhere (\S\ref{sec:results-controls}) and declines
here.

\begin{table}[!ht]\centering
\caption{Partner-regime table against a v0.5 opponent. Each row is the difference, in percentage
points of win-rate edge, between the edge of a team consisting of \sestina{} plus two copies of the
named partner and the edge of a team consisting of \vsix{} plus the same two partners, both against
three copies of the opponent; the interval is a between-cell contrast as in
\S\ref{sec:evaluation-intervals}. Per-bank values are given in bank order, 7090001 then 7090002.}
\label{tab:partners}
{\tabfont\input{tables_v07/partners_v05}}
\end{table}

Against a \vsix{} opponent all four rows are positive and clear zero
(Table~\ref{tab:partnerssix}), and this is where the measurement disagrees with the registered
expectation. The protocol states in advance \vsevenTrainOppSixSelf{} in self-play and
\vsevenTrainOppSixLo{} to \vsevenTrainOppSixHi{} under partner change. Measured:
\vsevenOppSixItself{} in self-play, and of the three changed rows two fall below the stated range,
namely \vsevenOppSixVthree{} against \texttt{v03}, more than a point below its floor, and
\vsevenOppSixDetective{} against \texttt{detective}.

\begin{table}[!ht]\centering
\caption{Partner-regime table against a \vsix{} opponent. Per-bank values are given in bank order,
7090001 then 7090002.}
\label{tab:partnerssix}
{\tabfont\input{tables_v07/partners_v06}}
\end{table}

\begin{table}[!ht]\centering
\caption{Transfer across partners: \sestina{} over \vsix{} against \vsix{} over v0.5, on the same
eight settings and the same banks. The two training rows are quoted from the registered
protocol; the protocol records that its v0.6-over-v0.5
training row is computed on three of seven changed settings, so that row is not directly comparable
with the holdout row above it. Reported and not gated.}
\label{tab:regimes}
{\tabfont\input{tables_v07/regimes}}
\end{table}

The protocol stated in advance that if the ordering in Table~\ref{tab:regimes} survived, the
appropriate summary is that \sestina{}'s advantage transfers across partners about as well as
\vsix{}'s did and, on the median-to-self-play ratio, slightly less well. The ordering survived
(\vsevenSoneRatio{} against \vsevenBaseRatio{}), so that is the paper's summary. The absolute
figures belong beside the ratio: \sestina{}'s changed-partner median is \vsevenSoneMedian{}
against \vsix{}'s \vsevenBaseMedian{}, and the ratio is smaller mainly because the self-play row it
divides by is more than six times larger.

A one-seat upgrade is worth \vsevenOneSeatFive{}~pp against a v0.5 opponent (registered expectation
\vsevenTrainOneSeatFive{}) and \vsevenOneSeatSix{}~pp against a \vsix{} opponent (registered
\vsevenTrainOneSeatSix{}). The single-seat upgrade produced a positive gain, which is consistent
with benefits beyond identical-agent coordination; because interacting team effects need not be
additive, the ratio of the one-seat gain to the three-seat gain is descriptive and does not
establish what share of the advantage is individual.

\subsection{Cross-play across training runs}
\label{sec:results-crossplay}

Diagonal mean \vsevenXpDiagonal{}, off-diagonal mean \vsevenXpOffDiagonal{}, off-diagonal minus
diagonal \vsevenXpGap{}~pp against a per-cell half-width of \vsevenXpHalfWidth{}; no interval is
computed for the off-diagonal-minus-diagonal contrast itself, so the comparison is descriptive
rather than a formal equivalence test. The worst off-diagonal cell (\vsevenXpWorstOff{}) is
\vsevenXpWorstOffShortfall{}~pp below the weakest diagonal cell (\vsevenXpWeakestDiag{}), and the
best off-diagonal (\vsevenXpBestOff{}) exceeds every diagonal cell. The registered claim that the
advantage is not a private convention between identical fits is satisfied on the point estimate
against a threshold of \vsevenStwoThreshold{}~pp. For scale, the coordination literature documents
convention collapses under cross-play that are large relative to self-play performance
\cite{bard-hanabi,hu-otherplay}, in score units that are not directly comparable with win-rate
points; no similarly large cross-play decrease was observed in these runs.
Table~\ref{tab:crossplay} is the cross-play matrix and Table~\ref{tab:crossh2h} the three pairs
head to head.

\begin{table}[!ht]\centering
\caption{Cross-play. Three runs of the same architecture, trained separately with the same fitting
infrastructure on disjoint fitting banks, with different search trajectories and, for one, a
different starting basin; pairwise parameter
distance \vsevenXpDistanceLo{} to \vsevenXpDistanceHi{} over \vsevenCoordinates{} coordinates.}
\label{tab:crossplay}
{\tabfont\input{tables_v07/crossplay}}
\end{table}

\begin{table}[!ht]\centering
\caption{The three separately trained runs head to head. The nominal intervals for
\vsevenHtwohSeparated{} of the three pairs exclude zero; this is an observed separation under the
reported procedure, not a coverage-validated test. That the three parameter vectors are distinct
(Table~\ref{tab:crossplay}) is a separate fact from a demonstrated performance difference.}
\label{tab:crossh2h}
{\tabfont\input{tables_v07/crossh2h}}
\end{table}

\subsection{Rule-variant sensitivity}
\label{sec:results-dialects}

Every row keeps its sign, every row replicates, and the largest excursion from the default row is
\vsevenDiaMaxExcursion{}~pp on \vsevenDiaMaxExcursionRow{}, inside the registered
\vsevenSfiveTolerance{}~pp tolerance. The advantage retains a positive estimate under 48-card
Literature, all three arbitration orders, both declaration-permission changes and the full legacy
variant. Table~\ref{tab:dialects} gives every row.

\begin{table}[!ht]\centering
\caption{The rule-variant table: \sestina{}'s win-rate edge in percentage points against \fcheap{}
under each variant, both configurations playing the same variant in every row, with nominal 95\% intervals
as in \S\ref{sec:evaluation-intervals}. Per-bank values are given in bank order, 7090003 then
7090001. This is the first sweep of its kind in the corpus and has no training-bank counterpart.}
\label{tab:dialects}
{\tabfont\begin{tabular}{lrrrr}
\toprule
Rule variant & Edge (pp) & 95\% CI & Per bank & vs default \\
\midrule
default & $+4.73$ & $[+4.10, +5.36]$ & $+4.57$ / $+4.88$ & $+0.00$ \\
no-out-of-turn & $+5.25$ & $[+4.63, +5.87]$ & $+4.72$ / $+5.78$ & $+0.52$ \\
no-cardless-declare & $+4.90$ & $[+4.27, +5.52]$ & $+4.71$ / $+5.08$ & $+0.17$ \\
maxasks=360 & $+4.73$ & $[+4.10, +5.36]$ & $+4.57$ / $+4.88$ & $+0.00$ \\
arb=high & $+4.85$ & $[+4.21, +5.49]$ & $+5.03$ / $+4.67$ & $+0.12$ \\
arb=turn & $+5.51$ & $[+4.88, +6.14]$ & $+5.26$ / $+5.77$ & $+0.78$ \\
sets=8 & $+4.13$ & $[+3.52, +4.73]$ & $+3.77$ / $+4.48$ & $-0.60$ \\
legacy & $+5.12$ & $[+4.50, +5.75]$ & $+4.62$ / $+5.63$ & $+0.40$ \\
\bottomrule
\end{tabular}
}
\end{table}

Two readings belong with the table. The \texttt{maxasks} row is bit-identical to the default
(\vsevenMaxasksIdentical{}), because the 400-ask cap is never reached in this population, so the cap
change produced no observed behavioural effect in this sample. Consequently the legacy residual,
which the protocol assigns to the forced-endgame willingness ladder for want of a flag of its own,
is arithmetically legacy minus two effective components rather than three; it is
\vsevenLegacyResidual{} $[\vsevenLegacyResidualLo, \vsevenLegacyResidualHi]$. That interval is a
quadrature propagation over cells that share banks and deals with unit coefficients, so it is not a
valid 95\% interval, and it would be wider than a paired difference if the relevant covariance were
non-negative (Appendix~\ref{app:corrections-self}); it is \vsevenLegacyResidualRatio{} times the size of the
residual. The residual estimate is not statistically interpretable under the current propagation,
and no separate ladder intervention was conducted.

\subsection{Adversarial-search results}
\label{sec:results-adversarial}

\vsevenBfourFits{} of \vsevenBfourArmsWord{} searches produced a weight vector and all
\vsevenBfourArmsWord{} were evaluated. None of the \vsevenBfourArmsWord{} searches found a positive
edge against \sestina{} at the tested budgets. Every one loses, by between
\vsevenBfourBestLoss{} and \vsevenBfourWorstLoss{}~pp; every sign replicates; and the largest upper
bound anywhere is \vsevenBfourWorstUpper{}, which is \vsevenBfourFloorGap{}~pp below the class
floor. The registered numerical criterion is satisfied using the nominal intervals.
Table~\ref{tab:adversaries} gives every search.

\begin{table}[!ht]\centering
\caption{\vsevenBfourArms{} separately fitted adversarial searches against \sestina{}, run with
common cross-entropy-method fitting infrastructure (the \texttt{tune} driver) and differing in
objective, initialisation and search variance. Each row is the adversarial response win-rate edge
in percentage points: the fitted adversary's edge
against \sestina{}, so a positive value would mean the adversary beats the target. Each was fitted
on a deal stream disjoint from the banks it is evaluated on. Per-bank values are given in bank
order, 7090001 then 7090002.}
\label{tab:adversaries}
{\tabfont\input{tables_v07/adversaries}}
\end{table}

The searches control all three opposing seats, run in the synchronised correlation regime, and use a
budget of \vsevenBfourGens{} generations of a \vsevenBfourPop{}-member population at
\vsevenBfourDeals{} deals per candidate. They maximise four different objectives, being win rate,
declaration error, event count and forced-endgame incidence, and two of them depart deliberately from
the template, one using a wide search variance and one seeded in the v0.5 basin, so that the set is
not one search repeated. The declaration-objective arm is read against the declaration-family floor of
\vsevenDeclFloor{} rather than the general floor.

These are observed response edges within the tested classes and budgets; they do not provide an
upper bound on best-response value. \S\ref{sec:discussion-exploit} states what would be required to
convert them into a statement about \sestina{}.

\subsubsection{Transcript-inversion adversary}
\label{sec:results-whitebox}

\texttt{Z08} is the adversary class the adversary model of \S\ref{sec:threat} was written for: an
opponent that reads a model of the target's deterministic policy and reconstructs a policy-aware
deal posterior from the public transcript, testing whether public policy specification among
identical deterministic teammates exposes exploitable behavioural information. It reaches
\vsevenZeightEdge{}~pp $[\vsevenZeightLo, \vsevenZeightHi]$, a replicated loss.

Three limits belong with that cell, since it is the one a sceptical reader should weigh most
heavily. The target model is a surrogate: the fixed transcript-inversion component in the base agent
queries a \vsix{} model of the target rather than the frozen configuration itself, a setting carried
forward from the phase whose target \emph{was} \vsix{}, so the likelihood is computed against a
policy that does not search, leaves the urgency escalation on, and carries neither the tie-break nor
the contestation coordinate; this is an inverter using a \vsix{} surrogate, not an exact
reconstruction of \sestina{}'s policy, and the corpus does not measure how far the surrogate's
choice diverges from the frozen target's, so the size of the resulting degradation is unquantified.
The inverter also accumulates per-card evidence additively across the game rather than carrying the
joint conditioning of \eqref{eq:coord} (\S\ref{sec:agent-coordination}). \vsevenZeightInert{} of its fitted coordinates are inert, because the white-box base derives from the
\vsix{} agent and reads only that policy's \vsevenVsixCoords{} coordinates while the search extends
its vector to \vsevenCoordinates{}; the effective search is therefore over \vsevenZeightLive{} live
coordinates plus the fixed transcript-inversion component in the base agent. And the budget is
small. The result is evidence that
this inverter, at this budget, does not exploit \sestina{}; it is not evidence that no white-box
attack can.

\paragraph{Training and evaluation discrepancy.} The in-class arm reaches
\vsevenZoneFitMargin{}~pp of paired margin in its final generation on its own fitting stream, and
\vsevenZoneFitPeak{}~pp at its best generation, before landing at \vsevenZoneEdge{}~pp on disjoint
banks. Most of that difference is not the change of bank: the same run re-scores its selected vector
on the fitting stream and already falls below parity there, because the generation figure is a
maximum over candidates at \vsevenBfourDeals{} deals each rather than a measurement of the vector
that was retained. Re-evaluation closes most of the gap and disjoint banks close the rest.

\subsection{Evaluation controls}
\label{sec:results-controls}

These controls assess selected failure modes of the evaluation procedure. Table~\ref{tab:controls}
is the planted-edge recovery.

\begin{table}[!ht]\centering
\caption{Planted-edge recovery, in percentage points of win-rate edge, pooled over both banks with
nominal 95\% intervals. \emph{Planted cost} is the handicap's cost to the target against \sestina{};
\emph{recovered excess} is the responder's edge against the handicapped target minus its edge
against the unhandicapped one, a between-cell contrast. The sub-floor rung must \emph{not} be
recovered.}
\label{tab:controls}
{\tabfont\input{tables_v07/controls}}
\end{table}

Two of the three supra-floor rungs are recovered with a lower bound above the
\vsevenPrimaryFloor{}~pp floor. The third, the \texttt{0.08} rung, recovers
\vsevenBnineHeightExcess{} $[\vsevenBnineHeightExcessLo, \vsevenBnineHeightExcessHi]$, above the floor
on the point estimate but with an interval straddling it, so on the lower-bound reading this paper
applies to its own primary claim it is not recovered. Its planted cost,
\vsevenBnineHeightCost{} $[\vsevenBnineHeightCostLo, \vsevenBnineHeightCostHi]$, straddles the floor
as well, so that rung is supra-floor only on a point estimate. Recovered size tracks planted size
monotonically across all four rungs.

The sub-floor rung is not recovered, at \vsevenBnineHfiveExcess{}
$[\vsevenBnineHfiveExcessLo, \vsevenBnineHfiveExcessHi]$: the sub-threshold control did not show a
statistically resolved recovered excess, which is the required outcome, although its upper bound
lies above the floor. Two qualifications belong with it: phase~2 measured the excess at this same handicap rung at
\vsevenSubFloorExcess{} $[\vsevenSubFloorExcessLo, \vsevenSubFloorExcessHi]$, so this cycle reaches
the same conclusion from a point estimate a little under half a point higher, and the upper bound
here grazes the floor. The protocol had named a different phase-2 row for this comparison;
Appendix~\ref{app:corrections-self} records the correction.

The identity control is exact: \sestina{} against itself is \vsevenIdentityRate{} with a
zero-width interval on both banks over \vsevenIdentityN{} games each. The responders fitted for these
controls also lose to unhandicapped \sestina{} in all four rows, a second sample of
\S\ref{sec:results-adversarial}'s finding from a different bank and against a handicapped target;
they maximise win rate as five of the \vsevenBfourArmsWord{} adversarial searches do, so this is a second sample
rather than an independent method. Table~\ref{tab:residual} is the seat-isolation residual.

\begin{table}[!ht]\centering
\caption{The seat-isolation residual on holdout material, \vsevenSsixCells{} cells over
\vsevenSsixDeals{} audited deals: decisions the S6 reconstruction could not reproduce from own hand,
public record, rules and reset seed, over decisions audited. The right column is the registered
zero-tolerance condition (one worker thread, agents rebuilt per deal); the left column is one
worker thread with reused agents. Neither column is the scored mode, which used
\vsevenScoredThreads{} worker threads with reused agents and was not audited on holdout.}
\label{tab:residual}
{\tabfont\input{tables_v07/residual}}
\end{table}

The gate column is zero for every configuration on both banks (\vsevenSsixGateNonzero{} nonzero
cells), so no inter-seat channel is detected under the registered test condition. The fresh-agent,
single-thread configuration passed the specified isolation tests. The scored configuration retained
a cross-deal state residual in searching agents, so these test results do not establish strict seat
isolation or exact reproducibility for the scored mode. The left column reproduces the known engine
defect of \S\ref{sec:agent-coordination} on holdout material and reproduces its shape:
\vsevenSsixPlainNonzero{} cells are nonzero and every one belongs to a \emph{searching}
configuration, with the non-searching blueprint zero on both banks. The audited decision counts
move between the two columns in \vsevenSsixDenomMoved{} of the cells, which is the measured form of
the statement in \S\ref{sec:evaluation-side} that rebuilding agents per deal changes play. The
residual's effect on a scored win rate is not bounded by a control cell; the largest discrepancy
observed between identical scored commands is one game in \vsevenResidueGames{}
(Appendix~\ref{app:corrections-reproducibility}), and a small observed residual does not establish
a bound on possible bias across all experiments. For the intervals of this paper the consequence is
the one stated in \S\ref{sec:evaluation-intervals}: the scored mode retains this state, so for
comparisons involving searching agents the independence of deal blocks and the coverage of the
intervals have not been established, and every interval is labelled nominal.

\subsection{Computational cost}
\label{sec:results-cost}

The record priced the frontier's search at \vsevenCostRecorded$\times$ the blueprint, a figure that
describes \fcheap{}. On a like-for-like basis the observed whole-match slowdown is
\vsevenCostFcheapTight$\times$ for \fcheap{} and \vsevenCostFrozenTight$\times$ for \sestina{},
measured against the cheapest opponent. All three figures exceed \vsevenCostRecorded$\times$, so the
record understates the cost. The whole-panel column is lower because expensive panel members place a
large common term on both sides of the ratio, and it is the most diluted of the three.
Table~\ref{tab:cost} gives the figures. The execution environment is recorded in
Appendix~\ref{app:repro}.

\begin{table}[!ht]\centering
\caption{Observed whole-match slowdown relative to the blueprint. The harness reports whole-match
throughput, which carries the opponent's cost on both sides of the ratio, so these ratios are
diluted by the opponent's cost and are closest to the agent-only cost against the cheapest
opponent; a strict lower bound on agent-only cost would require assuming that a policy does not
change game length or the opponent's work, which is not established. No figure carries an interval,
these being wall-clock timings taken across a battery that ran over two days, and the aggregation is
a ratio of medians, which differs from a median of per-bank ratios by at most \vsevenCostAggDelta{}
across the six median cells.}
\label{tab:cost}
{\tabfont\input{tables_v07/cost}}
\end{table}
